\documentclass[AER]{AEA}

% --- Packages ---
\usepackage{amsmath,amssymb,mathtools}
\usepackage{newtxtext}
\usepackage{newtxmath}
\usepackage{bm}
\usepackage{booktabs}
\usepackage{array}
\usepackage{placeins}
\usepackage{natbib}
\usepackage{enumitem}
\usepackage{url}
\usepackage[hidelinks]{hyperref}
\usepackage{xcolor}
\usepackage{tikz}

\urlstyle{same}

% 【修复】：给 blacksquare 加上数学模式符号 $ $
\newcommand{\qed}{\nolinebreak\hfill\ensuremath{\blacksquare}\par}

% --- 终极防伪与布尔开关 ---
\newcommand{\noopsort}[1]{} 

% 在正文全局中，让蓝三角隐身
\DeclareRobustCommand{\AERlink}[1]{}

% --- 定理环境 (AER标准：Proposition大写) ---
\newtheorem{proposition}{PROPOSITION}
\newtheorem{lemma}{LEMMA}
\newtheorem{theorem}{THEOREM}
\newtheorem{corollary}{COROLLARY}
\newtheorem{definition}{DEFINITION}
\newtheorem{remark}{REMARK}

% --- 核心设置 ---
\issueName{THE AMERICAN ECONOMIC RUBBISH}
\draftSpacing{1.5}
\setcounter{page}{23} % 接续上一篇页码
\pubMonth{April}
\pubYear{2026}
\pubVolume{1}
\pubIssue{2}

\begin{document}

% 1. 标题与十字架脚注
\title{Simping to Win: Contingent Courtship, Appropriable Attention, and Relational Contracts in the Affection Economy$^\dagger$}
\shortTitle{Simping to Win}

% 2. 作者与星号脚注
\author{Peter Price$^*$}
\date{\today}

\begin{abstract}
This paper develops a theory of ``simping'' as a commitment and appropriability problem in courtship. A suitor (the ``simp'') supplies costly attention and favors that can be partially appropriated even when the target does not commit. In a one-shot game, unconditional courtship that is fully appropriable cannot change partner choice. We introduce a \textit{contingency technology} that shifts courtship from appropriable transfers toward relationship-specific investments whose benefits are realized only under mutual commitment. We characterize the optimal mix of intensity and contingency, and provide a repeated-game foundation in which relational contracts endogenize contingency via continuation values and credible exit. The analysis highlights an efficiency rationale for platform features that facilitate transparent reciprocity.
\noindent (\textit{JEL} C72, C73, D86, D91, J12) \\
% \noindent \textbf{Keywords:} courtship, attention economy, reciprocity, commitment, relational contracts, dating markets
\end{abstract}

\maketitle

% 3. 脚注底层对齐黑魔法
\let\thefootnote\relax\footnotetext{\noindent\makebox[0em][r]{$^*$}Price: Department of Theoretical Romance, University of Affection (email: p.price@\allowbreak affection.edu). This is a theoretical paper studying incentive constraints and commitment in courtship. The model treats ``winning'' as the voluntary formation of a stable exclusive relationship. It does not advocate manipulatio or coercion; rather, it analyzes incentive compatibility and commitment problems.}
\let\thefootnote\relax\footnotetext{\noindent\makebox[0em][r]{$^\dagger$}Go to \url{https://webofnothing.top/noi/10.N0/aer.202604.A.0101.0004} to visit the article page for additional materials and author disclosure statements.}

% 【施法完毕】：恢复阿拉伯数字脚注
\def\thefootnote{\arabic{footnote}}
\setcounter{footnote}{0}

% 4. NOI 绝对定位纹身
\begin{tikzpicture}[remember picture, overlay]
  \node[anchor=north west, align=left, font=\scriptsize] at ([xshift=2cm, yshift=-1.5cm]current page.north west) {
    \textit{American Economic Rubbish 2026, 1(2): 23--33} \\
    \href{https://webofnothing.top/noi/10.N0/aer.202604.A.0101.0004}{\textit{https://webofnothing.top/noi/10.N0/aer.202604.A.0101.0004}}
  };
\end{tikzpicture}

% =========================================================
\section{Introduction}
% =========================================================

Courtship in modern ``attention economies'' often features a familiar pattern:
a high-investment suitor provides costly attention, favors, or gifts without receiving commitment in return.
In Chinese internet slang, such a suitor is called a \emph{ti\v{a}n g\v{o}u} (literally ``licking dog''; colloquially, a ``simp''),
a term referring to persistently one-sided romantic investment.\footnote{The term \emph{ti\v{a}n g\v{o}u} has become widely used in Chinese-language online communities to describe a person who invests heavily in a romantic interest without reciprocation.}

This paper provides a theoretical foundation for when such a suitor can \emph{win}---defined as the voluntary formation of a stable exclusive relationship---and when the strategy is doomed to be exploited.
Our central claim is simple but sharp: \emph{unconditional, fully appropriable courtship cannot change partner choice in a one-shot environment.}
To win, the suitor must redirect investment from ``appropriable attention'' toward \emph{contingent} and/or \emph{relationship-specific} actions whose benefits are realized primarily under mutual commitment.

The core friction is appropriability:
many courtship inputs (messages, gifts, errands, emotional validation) deliver immediate utility to the recipient regardless of future commitment.
If the recipient can accept these inputs and still match elsewhere, then unconditional generosity has little strategic leverage; it becomes a dominated instrument for persuasion.
This logic formalizes a widely observed intuition: ``giving more'' does not necessarily improve one's chance of being chosen if the other party can retain the benefit while pursuing superior outside options.

We introduce a reduced-form \emph{contingency index} capturing the share of courtship value that is \emph{not} appropriable without commitment.
Examples include:
(i) joint experiences requiring mutual participation,
(ii) investments that raise match quality only inside the relationship (e.g., learning the partner's preferences, integrating social networks), and
(iii) credible boundary-setting rules (withdrawal of attention after non-reciprocation).
We study the suitor's optimal choice of (intensity, contingency), characterize equilibrium relationship formation, and provide a repeated-game microfoundation in which contingency arises endogenously from relational contracting and credible exit.

\paragraph{Contributions.}
First, we model ``simping'' as a standard incentive problem combining (i) costly transfers and (ii) partial appropriability by a strategic recipient.
Second, we show in a one-shot model that fully appropriable gifts do not affect partner choice; winning requires contingent, relationship-specific instruments.
Third, we derive optimal contingency and show how high attachment helps only if it finances \emph{contingent} investment rather than unconditional transfers.
Fourth, we connect contingency to repeated-game continuation values and credibility: patience and outside options can substitute for formal contracting.
Finally, we discuss welfare and platform design: commitment and reciprocity tools can reduce inefficient overinvestment and mismatches.

\paragraph{Related literature.}
Our analysis draws on classic insights from signaling and commitment \citep{Spence1973, FudenbergTirole1991},
relational contracts \citep{BakerGibbonsMurphy2002, MailathSamuelson2006},
and matching and family economics \citep{Becker1973, GaleShapley1962}.
We interpret courtship inputs as noncontractible investments with partial appropriability, echoing hold-up logic in the theory of the firm, but focus on romantic partner choice rather than firm boundaries.

% =========================================================
\section{Environment and Definitions}
% =========================================================

We study a two-player game between a suitor~$S$ (the ``simp'') and a target~$T$ (the potential partner).

\subsection{Payoffs and Outside Options}

The target has an outside option $\bar{V} \in \mathbb{R}$ representing the utility from her best alternative (another match or remaining single).
The suitor obtains intrinsic utility $U \ge 0$ from forming an exclusive relationship with the target.

If $T$ forms a relationship with $S$, the target obtains baseline utility $b \in \mathbb{R}$ from match quality (chemistry and compatibility), plus the benefit from courtship inputs.
We intentionally allow $b$ to be below $\bar{V}$: the suitor may not be the target's ``default best'' match absent additional investment.

\subsection{Courtship Intensity and Contingency}

The suitor chooses two objects:

\begin{itemize}[leftmargin=0.75in]
\item \textbf{Intensity} $x \ge 0$: the total level of costly courtship input (time, effort, favors, resources).
\item \textbf{Contingency} $r \in [0,1]$: the fraction of the courtship value that is \emph{not appropriable} without commitment.
\end{itemize}

\noindent
\emph{Interpretation.}---The courtship input yields total value $x$ to the target \emph{if} the relationship forms.
If the target does \emph{not} commit, only a share $(1-r)x$ is nonetheless enjoyed (e.g., unilateral gifts or attention consumed),
while the remaining share $rx$ is effectively relationship-specific (e.g., joint activities, contingent effort, credible boundaries).

The suitor incurs cost $c(x)$ with $c(0)=0$, $c'(x)>0$, and $c''(x)\ge 0$.
Implementing contingency incurs additional (possibly psychological or social) cost $k(r)$ with $k(0)=0$, $k'(r)\ge 0$, and $k''(r)\ge 0$.

\begin{remark}[Why Contingency Is Costly]
In practice, making effort contingent can be socially awkward (e.g., initiating ``What are we?'' conversations), emotionally costly (e.g., withdrawing attention after non-reciprocation), or require commitment devices (e.g., accountability structures, cultivating alternative options).
These frictions are captured by $k(r)$.
\end{remark}

\subsection{Timing}

The game proceeds as follows:

\begin{enumerate}[leftmargin=0.75in]
\item $S$ chooses $(x,r)$.
\item $T$ observes $(x,r)$ and chooses $a \in \{C, R\}$, where $C$ denotes committing to an exclusive relationship with $S$ and $R$ denotes rejecting (i.e., matching elsewhere or remaining single).
\item Payoffs are realized.
\end{enumerate}

\subsection{Payoffs}

If $T$ commits ($a=C$):
\begin{align}
u_T(C) &= b + x, \label{eq:uTC}\\
u_S(C) &= U - c(x) - k(r). \label{eq:uSC}
\end{align}

If $T$ rejects ($a=R$):
\begin{align}
u_T(R) &= \bar{V} + (1-r)\,x, \label{eq:uTR}\\
u_S(R) &= -c(x) - k(r). \label{eq:uSR}
\end{align}

This formulation isolates the key asymmetry: the target can appropriate $(1-r)x$ even without committing.

\begin{definition}[Simping]
We say the suitor is ``simping'' when he chooses high $x$ paired with low $r$, so that most courtship value is appropriable upon rejection.
\end{definition}

\begin{definition}[Winning]
The suitor ``wins'' if the equilibrium outcome features $T$ choosing $C$ (exclusive relationship formation).
\end{definition}

% =========================================================
\section{Baseline: Unconditional Courtship Cannot Win}
% =========================================================

We first establish a stark benchmark: if courtship is fully appropriable upon rejection ($r=0$), then unconditional intensity cannot alter partner choice.

\begin{proposition}[Irrelevance of Fully Appropriable Courtship]\label{prop:irrelevance}
Suppose $r=0$ is the only feasible contingency level.
Then:
\begin{enumerate}[leftmargin=0.75in]
\item If $b < \bar{V}$, the target rejects ($R$) for any $x \ge 0$; the suitor cannot win.
\item In any subgame-perfect equilibrium, the suitor chooses $x=0$ (no courtship) and the target rejects whenever $b<\bar{V}$.
\end{enumerate}
\end{proposition}

Proposition~\ref{prop:irrelevance} captures the strategic core of why unconditional generosity fails in a one-shot environment:
when the recipient can keep what is given and still choose the better outside option, increasing gifts does not change the comparison.
The difference $u_T(R) - u_T(C) = \bar{V} - b$ is independent of $x$, so no amount of unconditional investment can alter the target's ranking.

\begin{remark}[Economics of ``Being Used'']
In the baseline, high $x$ with $r=0$ transfers surplus to $T$ without affecting $T$'s action.
Such behavior can be privately rational only if the suitor directly values the act of giving itself (not modeled here) or suffers from a commitment or self-control problem (discussed in Section~\ref{sec:extensions}).
\end{remark}

% =========================================================
\section{Contingent Courtship Design}
% =========================================================

We now allow $S$ to choose $r>0$ at cost $k(r)$, shifting courtship toward relationship-specific or contingent inputs.

\subsection{The Target's Best Response and the Commitment Constraint}

Given $(x,r)$, the target commits if and only if
\begin{equation}\label{eq:commitconstraint}
b + x \;\ge\; \bar{V} + (1-r)\,x
\quad \Longleftrightarrow \quad
b + r\,x \;\ge\; \bar{V}.
\end{equation}

Condition~\eqref{eq:commitconstraint} is the central incentive constraint: only the \emph{non-appropriable} portion $r\,x$ has persuasive power.
The term $r\,x$ represents the courtship value the target forfeits by rejecting, and it is this forgone value---not the total investment---that determines whether commitment is incentive compatible.

\subsection{Optimal Intensity Given Contingency}

Fix $r>0$ and suppose $b<\bar{V}$.
To induce commitment, the suitor chooses the minimal $x$ satisfying~\eqref{eq:commitconstraint}:
\begin{equation}\label{eq:xstar}
x(r) = \frac{\bar{V}-b}{r}.
\end{equation}
If $b\ge \bar{V}$, then $x(r)=0$ suffices: the suitor wins without costly courtship.

\subsection{Optimal Contingency}

Assume $b<\bar{V}$ and $r\in(0,1]$ is feasible.
Substituting~\eqref{eq:xstar} into $u_S(C)$ yields the suitor's problem:
\begin{equation}\label{eq:rsproblem}
\max_{r\in(0,1]} \;
U - c\!\left(\frac{\bar{V}-b}{r}\right) - k(r).
\end{equation}

\begin{proposition}[Characterization of Optimal Contingency]\label{prop:optcontingency}
Suppose $c$ and $k$ are differentiable with $c'>0$ and $k'\ge 0$.
If an interior solution $r^*\in(0,1)$ exists, it satisfies the first-order condition
\begin{equation}\label{eq:foc}
\frac{\bar{V}-b}{(r^*)^2}\;c'\!\left(\frac{\bar{V}-b}{r^*}\right) = k'(r^*).
\end{equation}
Moreover, the suitor wins (i.e., chooses $(x,r)$ that induces $C$) if and only if
\begin{equation}\label{eq:participation}
U - c\!\left(\frac{\bar{V}-b}{r^*}\right) - k(r^*) \;\ge\; 0.
\end{equation}
\end{proposition}

Equation~\eqref{eq:foc} equates the marginal benefit of raising contingency (reducing the required intensity) to the marginal cost of implementing contingency.
Condition~\eqref{eq:participation} is the suitor's participation constraint: winning must yield non-negative surplus.

\subsection{Closed-Form Example}

To make the comparative statics transparent, consider the following functional forms:

\[
c(x)=x
\quad\text{and}\quad
k(r)=\frac{\kappa}{2}\,r^2
\quad (\kappa>0).
\]

Then, for $b<\bar{V}$, the suitor's objective becomes
\[
U - \frac{\bar{V}-b}{r} - \frac{\kappa}{2}\,r^2.
\]
The interior optimum solves $(\bar{V}-b)/r^2=\kappa\, r$, yielding
\begin{equation}\label{eq:rclosed}
r^*=\min\left\{1,\;\left(\frac{\bar{V}-b}{\kappa}\right)^{1/3}\right\},
\qquad
x^*=\frac{\bar{V}-b}{r^*}.
\end{equation}

\begin{corollary}[Comparative Statics]\label{cor:comparative}
In the closed-form example above, holding $U$ fixed:
\begin{enumerate}[leftmargin=0.75in]
\item A larger competition gap $(\bar{V}-b)$ increases $r^*$ (weakly) and increases the required intensity $x^*$.
\item A higher contingency cost $\kappa$ decreases $r^*$ and increases $x^*$.
\item The suitor wins if and only if $U \ge (\bar{V}-b)/r^* + (\kappa/2)(r^*)^2$.
\end{enumerate}
\end{corollary}

\begin{remark}[How the Simp Can Win]
A high-attachment suitor (large $U$) can finance the required investment, but \emph{only} if he shifts that investment toward high-$r$ instruments.
If his investment remains fully appropriable ($r\approx 0$), the target's incentive constraint cannot be satisfied unless $b\ge \bar{V}$.
Thus, ``winning'' requires transforming unconditional giving into contingent, relationship-specific value creation.
\end{remark}

% =========================================================
\section{A Repeated-Game Foundation: Relational Contracts and Credible Exit}
% =========================================================

The reduced-form parameter $r$ can be microfounded by dynamic incentives.
We consider an infinite-horizon repeated interaction in which exclusivity is sustained via continuation values rather than enforceable contracts.

\subsection{Stage Game}

Time is indexed by $t=0,1,2,\dots$.
At the start of each period, if the relationship is ``active,'' the suitor chooses effort $x\ge 0$ (for simplicity, held constant over time).
The target then chooses whether to remain exclusive (cooperate) or to deviate (pursue outside options).

Per-period payoffs are as follows:
\begin{itemize}[leftmargin=0.75in]
\item \emph{Cooperation:} $T$ receives $b+x$ and $S$ receives $U - c(x)$.
\item \emph{Deviation:} $T$ receives $\bar{V}+x$ in the deviation period (she retains the contemporaneous appropriable benefit), and the relationship may terminate depending on the suitor's response.
\end{itemize}

\noindent Both players discount future payoffs by $\delta\in(0,1)$.

\subsection{Relational Enforcement and Effective Contingency}

Consider a grim-trigger strategy profile:
the suitor provides effort $x$ as long as the target has never deviated; after any deviation, the suitor permanently exits and effort drops to zero.
The target cooperates as long as the suitor provides $x$; otherwise she takes the outside option.

The target's one-shot gain from deviating is $\bar{V}-b$ (since she enjoys the current $x$ regardless).
The future loss from deviating is the discounted stream of relationship surplus $(b+x)-\bar{V}$.

\begin{proposition}[Sustaining Exclusivity via Continuation Values]\label{prop:repeated}
Under grim-trigger strategies, exclusivity is incentive compatible for the target if and only if
\begin{equation}\label{eq:repeatedIC}
b + \delta\, x \;\ge\; \bar{V}.
\end{equation}
\end{proposition}

Condition~\eqref{eq:repeatedIC} mirrors the one-shot constraint $b+r\,x\ge \bar{V}$ with \emph{effective contingency} $r=\delta$.
Intuitively, only the continuation value (scaled by $\delta$) is non-appropriable: the target can appropriate current-period effort by deviating after receiving it, but forfeits all future relationship surplus.

\subsection{Credibility of Exit: When the Simp Cannot Punish}

Relational contracts also require that punishment be credible.
If the suitor is so attached that he cannot credibly exit after deviation, the target anticipates no future loss and will deviate with impunity.

We model credibility by endowing the suitor with an outside option $W\ge 0$ available after exit.
If he stays after deviation, suppose his continuation payoff is $\underline{U}$ (e.g., reflecting jealousy, distrust, or diminished relationship quality).
Exit is credible if and only if $W \ge \underline{U}$.

\begin{proposition}[Credible Exit as a Discipline Device]\label{prop:credibility}
Suppose that after a deviation the suitor chooses whether to exit (receiving $W$ in perpetuity) or stay (receiving $\underline{U}$ in perpetuity), with $W$ and $\underline{U}$ common knowledge.
Then:
\begin{enumerate}[leftmargin=0.75in]
\item If $W < \underline{U}$, the suitor's threat to exit is not credible; in any Markov-perfect equilibrium the target deviates and exclusivity cannot be sustained unless $b\ge \bar{V}$.
\item If $W \ge \underline{U}$, exit is credible; exclusivity can be sustained whenever condition~\eqref{eq:repeatedIC} holds.
\end{enumerate}
\end{proposition}

\begin{remark}[The Strategic Self-Improvement Channel]
Proposition~\ref{prop:credibility} formalizes a common intuition: a suitor who lacks viable outside options cannot credibly enforce reciprocity.
In the model, increasing $W$ (through social networks, alternative matches, self-respect, or higher opportunity cost of time) strengthens the relational contract by making exit credible.
This does not require hostility or manipulation; it is simply the credibility condition for any self-enforcing agreement.
\end{remark}

% =========================================================
\section{Extensions}\label{sec:extensions}
% =========================================================

\subsection{Competition Among Multiple Suitors}

Let there be $n\ge 1$ alternative suitors generating the target's outside option $\bar{V}(n)$, increasing in $n$.
Then the competition gap $\bar{V}(n)-b$ rises with $n$, tightening both the one-shot constraint $b+r\,x\ge \bar{V}(n)$ and the repeated-game condition $b+\delta\, x\ge \bar{V}(n)$.
Thus, greater competition requires either (i)~higher contingency (stronger reciprocity norms or boundaries), (ii)~greater intensity, or (iii)~higher patience to magnify continuation values.

\subsection{Heterogeneity in Reciprocity Norms}

One interpretation of $r$ is technological: it measures what fraction of courtship value is inherently relationship-specific.
Another interpretation is behavioral: some targets experience a psychological cost from appropriating attention without reciprocation, effectively raising $r$ endogenously.
If such ``reciprocity norms'' vary across targets, then the same suitor action may be persuasive for one target but strategically irrelevant for another.
This heterogeneity predicts selection effects: unconditional giving appears to ``work'' in subpopulations with high endogenous~$r$, but fails systematically where appropriation is socially or psychologically inexpensive.

\subsection{Private Information and Screening}

Suppose the target's outside option $\bar{V}$ is private information, drawn from a distribution $F$ known to the suitor.
The suitor can offer a menu of courtship packages $\{(x_i,r_i)\}_{i}$ (e.g., a sequence of escalating contingent investments).
Under standard single-crossing conditions, an optimal menu is monotone: targets with higher outside options require higher effective contingent value $r\,x$ to commit.
In practice, this corresponds to \emph{gradual escalation of relationship-specific investments} rather than large unconditional gifts offered early, thereby reducing losses when $\bar{V}$ turns out to be high.

\subsection{Platform Design: Escrow and Verifiable Reciprocity}

In many markets, escrow mechanisms transform noncontractible exchange into contractible exchange.
Analogously, dating platforms could implement verifiable reciprocity features---for example, mutual exclusivity flags, transparent effort commitments, or time-bounded escalation protocols.
In the model, such tools effectively lower $k(r)$ (the cost of implementing contingency) and raise the feasible level of~$r$, thereby expanding the parameter region in which mutually beneficial relationships form and reducing wasteful appropriable transfers.

% =========================================================
\section{Welfare and Discussion}
% =========================================================

We briefly discuss efficiency implications.
Let total surplus from an exclusive relationship (absent courtship effort) be $U+b$, and total surplus from the outside option be $0+\bar{V}$ (normalizing the suitor's outside payoff to zero).
The first-best matching rule is: form the relationship if and only if $U+b \ge \bar{V}$.

However, Proposition~\ref{prop:irrelevance} shows that even when $U+b \ge \bar{V}$, the relationship may fail to form if the target's private incentive constraint is not satisfied.
This constitutes an inefficiency driven by appropriability and the absence of enforceable commitment.
Contingency tools (higher $r$) and relational contracting (higher $\delta$ coupled with credible exit) move outcomes closer to first best by aligning the target's private incentives with social surplus.

\begin{remark}[A Normative Reading]
The model does \emph{not} prescribe ``invest more to win.''
Rather, it yields three prescriptions:
(i)~unconditional giving is strategically weak under appropriability;
(ii)~winning requires shifting effort toward mutually meaningful, relationship-specific value and clear reciprocity; and
(iii)~credible exit (i.e., maintaining boundaries) is central to any self-enforcing relationship contract.
\end{remark}

% =========================================================
\section{Conclusion}
% =========================================================

This paper modeled ``simping'' as an appropriability and commitment problem in courtship.
In a one-shot setting, fully appropriable gifts cannot change partner choice: the recipient can accept them and still pursue the superior outside option.
Winning requires contingency---courtship inputs whose value is realized primarily under mutual commitment.
We characterized optimal intensity and contingency, showed how competition and contingency costs shape equilibrium outcomes, and provided a repeated-game foundation in which patience and credible exit endogenize effective contingency.

A key implication is conceptual: a high-attachment suitor can prevail not by increasing unconditional transfers, but by reallocating effort toward relationship-specific value and credible reciprocity.
This reallocation is both strategically necessary (to satisfy the target's incentive-compatibility constraint) and normatively preferable (to avoid inefficient overinvestment and exploitation).

Several directions merit future investigation.
First, the model could be enriched with two-sided investment, where both parties make relationship-specific contributions.
Second, introducing learning about match quality would allow the analysis to address the optimal timing and pacing of courtship escalation.
Third, empirical work using data from dating platforms could test the comparative-statics predictions regarding contingency, competition, and relationship formation rates.

\FloatBarrier

% =========================================================
% REFERENCES
% =========================================================

% 【唤醒蓝三角】：此时 DeclareRobustCommand 获得生命！
\DeclareRobustCommand{\AERlink}[1]{\makebox[0pt][r]{\href{#1}{\color{blue}$\blacktriangleright$}\hspace{0.3em}}}

% \nocite{*}
\bibliographystyle{aea}
\bibliography{AERub_Simp}

\appendix
\section{Appendix: Proofs}

% \subsection{Proof of Proposition \ref{prop:irrelevance}}
% Fix $r=0$. For any $x \ge 0$, $u_T(R) - u_T(C) = \bar{V} - b$. If $b < \bar{V}$, rejection is a dominant strategy. Anticipating this, $S$ chooses $x=0$ to minimize costs. \qed

% \subsection{Proof of Proposition \ref{prop:optcontingency}}
% The suitor chooses minimal $x = (\bar{V}-b)/r$ to induce $C$. Substituting into the objective and taking FOC with respect to $r$ yields the result via the chain rule. \qed

\subsection{Proof of Proposition \ref{prop:irrelevance}}

Fix $r=0$.
For any $x\ge 0$, the target's payoff from committing is $b+x$ and from rejecting is $\bar{V}+(1-0)x=\bar{V}+x$.
Thus $u_T(R)-u_T(C)=\bar{V}-b$, which is independent of $x$.
If $b<\bar{V}$, then $u_T(R)>u_T(C)$ for all $x$, so the target rejects regardless of $x$.
Anticipating rejection, the suitor's payoff is $u_S(R)=-c(x)-k(0)=-c(x)$, which is maximized at $x=0$ since $c$ is increasing and $c(0)=0$.
\hfill $\blacksquare$

\subsection{Proof of Proposition \ref{prop:optcontingency}}

Assume $b<\bar{V}$.
For any $r>0$, the target commits if and only if $x\ge (\bar{V}-b)/r$ by condition~\eqref{eq:commitconstraint}.
Since the suitor's payoff under commitment is $U-c(x)-k(r)$ and $c$ is increasing, the suitor optimally chooses the minimal $x(r)=(\bar{V}-b)/r$.
Substituting yields the optimization problem~\eqref{eq:rsproblem}.
If the optimum is interior, the first-order condition is
\[
-\frac{d}{dr}\,c\!\left(\frac{\bar{V}-b}{r}\right) - k'(r) = 0.
\]
By the chain rule,
\[
\frac{d}{dr}\,c\!\left(\frac{\bar{V}-b}{r}\right)
=
c'\!\left(\frac{\bar{V}-b}{r}\right)\cdot\left(-\frac{\bar{V}-b}{r^2}\right).
\]
Substituting and rearranging gives equation~\eqref{eq:foc}.
Finally, the suitor chooses to induce commitment (rather than abstaining) only if his equilibrium payoff is non-negative, yielding condition~\eqref{eq:participation}.
\hfill $\blacksquare$

\subsection{Proof of Proposition \ref{prop:repeated}}

Under perpetual cooperation, the target's discounted payoff is
\[
V_C = \sum_{t=0}^{\infty}\delta^t(b+x)=\frac{b+x}{1-\delta}.
\]
If the target deviates in the current period, she obtains $\bar{V}+x$ now (appropriating current effort) and then receives $\bar{V}$ in every subsequent period after the suitor exits:
\[
V_D = (\bar{V}+x)+\sum_{t=1}^{\infty}\delta^t\,\bar{V} = (\bar{V}+x)+\frac{\delta\,\bar{V}}{1-\delta}.
\]
Cooperation is incentive compatible if and only if $V_C\ge V_D$.
Multiplying both sides by $(1-\delta)$ yields
\[
b+x \;\ge\; (1-\delta)(\bar{V}+x)+\delta\,\bar{V} \;=\; \bar{V} + (1-\delta)\,x,
\]
which simplifies to $b+\delta\, x\ge \bar{V}$.
\hfill $\blacksquare$

\subsection{Proof of Proposition \ref{prop:credibility}}

\emph{Part 1.}
If $W<\underline{U}$, then after any deviation the suitor strictly prefers staying to exiting, so the exit threat is not sequentially rational.
Anticipating this, the target expects no loss of future relationship surplus from deviating; in the stage game she prefers $\bar{V}+x$ to $b+x$ whenever $b<\bar{V}$, so deviation occurs in every period.

\emph{Part 2.}
If $W\ge \underline{U}$, then exit is sequentially rational after deviation.
Given that the suitor will exit following any deviation, the target's incentive constraint reduces to the standard relational condition derived in Proposition~\ref{prop:repeated}.
\hfill $\blacksquare$

\end{document}